\section{USAJMO 2024}
        \begin{enumerate}
        	\item (\textit{USAJMO 2024}) Let $ABCD$ be a cyclic quadrilateral with $AB = 7$ and $CD = 8$. Points $P$ and $Q$ are selected on segment $AB$ such that $AP = BQ = 3$. Points $R$ and $S$ are selected on segment $CD$ such that $CR = DS = 2$. Prove that $PQRS$ is a cyclic quadrilateral.


 



\textbf{Solution 1}

If $\overline {AB}\parallel \overline{CD}$, then $PQRS$ is an isosceles trapezoid, and thus is cyclic. Otherwise, let $X = \overleftrightarrow{AB} \cap \overleftrightarrow{CD}$, WLOG with $A, C$ on segments $\overline{XB}, \overline{XD}$. By Power of a Point, 
 \[XA \cdot XB=XC \cdot XD \implies (XA)(XA+7)=(XC)(XC+8) \implies XA^2+7XA=XC^2+8XC\] 
 \[\implies XA^2+7XA+12=XC^2+8XC+12 \implies (XA+3)(XA+4)=(XC+2)(XC+6)\implies XP\cdot XQ = XR\cdot XS\] as desired $\square$

 -sillybone

 



\textbf{Solution 2}

First, let $E$ and $F$ be the midpoints of $AB$ and $CD$, respectively. It is clear that $AE=BE=3.5$, $PE=QE=0.5$, $DF=CF=4$, and $SF=RF=2$. Also, let $O$ be the circumcenter of $ABCD$. 

 
                    \begin{figure}[H]
                        \centering
                        \includegraphics[width=0.5\textwidth]{Solution/USAJMO_I_2024_1_Solution_Figure_1.png}
                    \end{figure}
                    

 By properties of cyclic quadrilaterals, we know that the circumcenter of a cyclic quadrilateral is the intersection of its sides' perpendicular bisectors. This implies that $OE\perp AB$ and $OF\perp CD$. Since $E$ and $F$ are also bisectors of $PQ$ and $RS$, respectively, if $PQRS$ is indeed a cyclic quadrilateral, then its circumcenter is also at $O$. Thus, it suffices to show that $OP=OQ=OR=OS$. 

 Notice that $PE=QE$, $EO=EO$, and $\angle QEO=\angle PEO=90^\circ$. By SAS congruency, $\Delta QOE\cong\Delta POE\implies QO=PO$. Similarly, we find that $\Delta SOF\cong\Delta ROF$ and $OS=OR$. We now need only to show that these two pairs are equal to each other. 

 Draw the segments connecting $O$ to $B$, $Q$, $C$, and $R$. 

 
                    \begin{figure}[H]
                        \centering
                        \includegraphics[width=0.5\textwidth]{Solution/USAJMO_I_2024_1_Solution_Figure_2.png}
                    \end{figure}
                    

 Also, let $r$ be the circumradius of $ABCD$. This means that $AO=BO=CO=DO=r$. Recall that $\angle BEO=90^\circ$ and $\angle CFO=90^\circ$. Notice the several right triangles in our figure. 

 Let us apply Pythagorean Theorem on $\Delta BEO$. We can see that $EO^2+EB^2=BO^2\implies EO^2+3.5^2=r^2\implies EO=\sqrt{r^2-12.25}.$ 

 Let us again apply Pythagorean Theorem on $\Delta QEO$. We can see that $QE^2+EO^2=QO^2\implies0.5^2+r^2-12.25=QO^2\implies QO=\sqrt{r^2-12}.$ 

 Let us apply Pythagorean Theorem on $\Delta CFO$. We get $CF^2+OF^2=OC^2\implies4^2+OF^2=r^2\implies OF=\sqrt{r^2-16}$. 

 We finally apply Pythagorean Theorem on $\Delta RFO$. This becomes $OF^2+FR^2=OR^2\implies r^2-16+2^2=OR^2\implies OR=\sqrt{r^2-12}$. 

 This is the same expression as we got for $QO$. Thus, $OQ=OR$, and recalling that $OQ=OP$ and $OR=OS$, we have shown that $OP=OQ=OR=OS$, and we are done.

 



\textbf{Solution 3}

We can consider two cases: $AB \parallel CD$ or $AB \nparallel CD.$ The first case is trivial, as $PQ \parallel RS$ and we are done due to symmetry. For the second case, WLOG, assume that $A$ and $C$ are located on $XB$ and $XD$ respectively. Extend $AB$ and $CD$ to a point $X,$ and by Power of a Point, we have 
 \[XA\cdot XB = XC \cdot XD,\] which may be written as 
 \[XA \cdot (XA+7) = XC \cdot (XC+8),\] or 
 \[XA^2 + 7XA = XC^2 + 8XC.\] We can translate this to 
 \[XA^2 + 7XA +12 = XC^2 + 8XC +12,\] so 
 \[XP\cdot XQ = (XA+3)(XA+4)=(XC+2)(XC+6)= XR\cdot XS,\] and therefore by the Converse of Power of a Point $PQRS$ is cyclic, and we are done.

 - \href{https://artofproblemsolving.com/wiki/index.php?title=User:Spectraldragon8}{spectraldragon8}

 



\textbf{Solution 4}

All 4 corners of $PQRS$ have equal power of a point ($12$) with respect to the circle $(ABCD)$, with center $O$.

 Draw diameters (of length $AQ$) of circle $(ABCD)$ through $Q$ and $S$, with length $A$.  Let $q$ be the distance from $Q$ to the circle along a diameter, and  likewise $s$ be distance from $S$ to the circle.

 Then $q(AQ-q) = s(AQ-s) = 12$ and $q,s < AQ/2$ (radius). Therefore, $q=s$ and $AQ/2 -q = AQ/2 -s$. But $AQ-q=OQ$, $AQ-s=OS$, $OQ = OP$ and $OS = OR$ by symmetry around the perpendicular bisectors of $PQ$ and $RS$, so $P,Q,R,S$ are all equidistant from $O$, forming a circumcircle around $PQRS$.

 -BraveCobra22aops, oinava

 


 



\textbf{Solution 4 (Coord Bash)}

Let $A(2a_1,2a_2)$, $B(2b_1,2b_2)$, $C(2c_1,2c_2)$, $D(2d_1,2d_2)$, and the circumcenter of quadrilateral $ABCD$ be $O(0,0)$.

 Let's list what we know from the givens:Since the radii of a circle are equal in length, we can let 
 \[k=a_1^2+a_2^2=b_1^2+b_2^2=c_1^2+c_2^2=d_1^2+d_2^2.\]

 From the distance formula on $AB$ and $CD$, we can simplify and get 
 \[a_1^2+a_2^2+b_1^2+b_2^2-2(a_1b_1+a_2b_2)=\frac{49}{4} \Rightarrow 2(a_1b_1+a_2b_2)=2k-\frac{49}{4}\] and 
 \[c_1^2+c_2^2+d_1^2+d_2^2-2(c_1d_1+c_2d_2)=16 \Rightarrow 2(c_1d_1+c_2d_2)=2k-16\] using the above substitution with $k$.

 Since $P$, $Q$, $R$, and $S$ are weighted points on $AB$ and $CD$, we can get 
 \[P\left(\frac{8a_1+6b_1}{7},\frac{8a_2+6b_2}{7}\right) \ \text{and} \ Q\left(\frac{6a_1+8b_1}{7},\frac{6a_2+8b_2}{7}\right)\] along with 
 \[R\left(\frac{3c_1+d_1}{2},\frac{3c_2+d_2}{2}\right) \ \text{and} \ S\left(\frac{c_1+3d_1}{2},\frac{c_1+3d_1}{2}\right).\]

 Now, we want to show that the perpendicular bisectors of the sides of quadrilateral $PQRS$ are concurrent to prove that it is cyclic. Moreover, we suspect that the circumcenter of quadrilateral $PQRS$ is the circumcenter of quadrilateral $ABCD$, both of them being $(0,0)$.

 We only list the linear equations of the perpendicular bisectors of lines $PS$ and $QR$, as the perpendicular bisectors of lines $PQ$ and $RS$ are the same as the perpendicular bisectors of lines $AB$ and $CD$, respectively.

 First, we analyze the perpendicular bisector of line $PS$. Notice that it has slope 
 \[-\frac{16a_1+12b_1-7c_1-21d_1}{16a_2+12b_2-7c_2-21d_2}\] and point 
 \[\left(\frac{16a_1+12b_1+7c_1+21d_1}{28}, \frac{16a_2+12b_2+7c_2+21d_2}{28}\right).\] The point-slope form equation would thus be 
 \[y-\frac{16a_2+12b_2+7c_2+21d_2}{28}=-\frac{16a_1+12b_1-7c_1-21d_1}{16a_2+12b_2-7c_2-21d_2}\left(x-\frac{16a_1+12b_1+7c_1+21d_1}{28}\right).\] Since we claim that $(0,0)$ is on the line, we substitute to get 
 \[(16a_1+12b_1+7c_1+21d_1)(16a_1+12b_1-7c_1-21d_1)=-(16a_2+12b_2+7c_2+21d_2)(16a_2+12b_2-7c_2-21d_2),\] and this simplifies to 
 \[16\left(16(a_1^2+a_2^2)+24(a_1b_1+a_2b_2)+9(b_1^2+b_2^2)\right)=49\left((c_1^2+c_2^2)+6(c_1d_1+c_2d_2)+9(d_1^2+d_2^2)\right).\] Using substitutions from the first three equations, this becomes 
 \[16\left(16k+12\left(2k-\frac{49}{4}\right)+9k\right)=49\left(k+3(2k-16)+9k\right)\] 
 \[\Rightarrow 16(49k-3\cdot49)=49(16k-3\cdot16)\] 
 \[\Rightarrow k-3=k-3,\] which is true, implying that $(0,0)$ does indeed satisfy the equation.

 Analogously, we can show that the perpendicular bisector of line $QR$ also passes through the origin.

 Since the perpendicular bisectors of the sides of quadrilateral $PQRS$ all intersect at the same point, namely $(0,0)$, which is also the circumcenter of quadrilateral $ABCD$, we can conclude that $PQRS$ is a cyclic quadrilateral.

 



\textbf{Solution 5 (Euclidean)}

Let $O$ be the center of the circumcircle of $ABCD$, and let $r$ be the circumradius of that quadrilateral. Also, let $E$ be the midpoint of $\overline{AB}$ and let $F$ be the midpoint of $\overline{CD}$. If we look at $\triangle AEO$, we see that it is a right triangle with hypotenuse $r$ and one leg $\frac{7}{2}$. Therefore, the other leg, $\overline{EO}$, has length $\sqrt{r^{2}-\frac{49}{4}}$. We have that $EP = EQ = \frac{1}{2}$, and therefore, by the Pythagorean Theorem we have that $OP = OQ = \sqrt{r^{2}-12}$.

 Applying the same logic to $\overline{CD}$, we find that $OS = OR = \sqrt{r^{2}-12}$. Because $PO = QO = SO = RO$, the points $P$, $Q$, $S$, and $R$ all lie on a common circle with center $O$ and radius $\sqrt{r^{2}-12}$. Therefore, $PQRS$ is a cyclic quadrilateral.

 


	\item (\textit{USAJMO 2024}) Let $m$ and $n$ be positive integers. Let $S$ be the set of integer points $(x,y)$ with $1\leq x\leq2m$ and $1\leq y\leq2n$. A configuration of $mn$ rectangles is called \textit{happy} if each point in $S$ is a vertex of exactly one rectangle, and all rectangles have sides parallel to the coordinate axes. Prove that the number of \textit{happy} configurations is odd.


 



\textbf{Solution 1}

We will proceed by induction on $n$. For the purpose of notation, we will group the points in $S$ into $m$ rows and $n$ columns and refer to them as such accordingly.

 
\textbf{Base Case}

 Consider the case where $n=1$. It is clear that any rectangle in the happy configuration is determined only by which two rows its vertices lie on, and that these two rows are automatically filled by the vertices of the rectangle. By repeatedly picking two rows to form a new rectangle from the remaining rows, the number of rectangles in this case is

 
 \[\frac{1}{m!}\prod_{k=1}^{m}\binom{2k}{2}=\frac{1}{m!}\binom{2m}{2}\binom{2m-2}{2}\cdots\binom{2}{2}=\frac{(2m)!}{m!\cdot2^m}\]

 We can evaluate the parity of this expression by expanding the factorials:

 
 \[\frac{(2m)(2m-1)(2m-2)\cdots(2)(1)}{(2m)(2m-2)\cdots(2)}=(2m-1)(2m-3)\cdots(1)\]

 Since none of the factors are even, this expression must be odd. As a result, there are always an odd number of happy configurations if $n=1$.

 
\textbf{Inductive Step}

 Assume that there are an odd number of happy configurations for some $m$ and $n$. We will now prove that there must also be an odd number of happy configurations for $m$ and $n+1$.

 Consider $S$ for $m$ and $n+1$. Denote the rightmost two columns of $S$ as the "strip".

 
Case 1: All rectangles in the happy configuration lie entirely within the strip or entirely outside of the strip.

 By the inductive hypothesis, there are an odd number of happy configurations in the region of $S$ outside of the strip. Additionally, by the base case, there are an odd number of happy configurations in the strip. Since no rectangles extend between the two regions, the number of happy configurations in this case is just the product of the two values, which is clearly odd.

 
Case 2: There exists a rectangle $R$ which contains some vertices inside of the strip and some vertices outside of the strip.

 Let a happy configuration with this property be denoted as $H$.

 First of all, because each vertex in a rectangle shares a column with another vertex and because the strip is only two columns wide, $R$ must have exactly two vertices in the strip. Furthermore, the two vertices in the strip must lie in the same column, and the same applies for the two vertices outside of the strip.

 Notice that both vertices of $R$ must lie on rows which each contain only one other point inside of the strip. This forces the existence of another rectangle $R'$ which has exactly two vertices in the strip and has vertices lying on the same two rows as $R$. Then, it must be possible to create a new happy configuration $H'$ by swapping the vertices of $R$ and $R_1$ which in the strip between the two rectangles.

 Notice that repeating this procedure on the new happy configuration results in $H$ again. As a result, there exists a one-to-one pairing between happy configurations for which such a rectangle $R$ exists. 

 However, the existence of such a pairing implies that the number of happy configurations in this case is even.

 
Conclusion:

 The number of happy configurations for $m$ and $n+1$ is the sum of the number of happy configurations from each of the two cases. Since there are an odd number of happy configurations in case 1 and an even number of happy configurations in case 2, there must be an odd total number of happy configurations.

 By induction, there must be an odd number of happy configurations of $S$ for any positive integer $m$ and $n$.

 
Note: It is possible to reduce the base case to $m=n=1$ because the inductive step can be applied to incrementing $m$ as well.

 
- pomme\_de\_terre\_ (aka BlackDaddyUWU)

 


 



\textbf{Solution 2}

We can group configurations into pairs where:

 
            \begin{enumerate}[label=(\roman*)]
                \item 1. If it is \textit{not} horizontally symmetric, it makes a pair with its horizontal reflection.
            \end{enumerate}
            
            \begin{enumerate}[label=(\roman*)]
                \item 2. Else, if it is \textit{not} vertically symmetric, it makes a pair with its vertical reflection.
            \end{enumerate}
            
            \begin{enumerate}[label=(\roman*)]
                \item 3. Else, if it is \textit{not} half-symmetric, we can similarly make a pair with its horizontal or vertical half-reflection.
            \end{enumerate}
            The \textbf{half-reflection} is defined by reflecting only the edge of each rectangle closest to the horizontal/vertical center, while keeping the other edge fixed. For example, a vertical half-reflection would send the rectangle
 \[\{(0, 0), (1, 0), (1, n-1), (0, n-1)\}\]to
 \[\{(0, 0), (1, 0), (1, n+1), (0, n+1)\}.\]Since every configuration that is not symmetric or half-symmetric comes in pairs, we only need to look at the parity of the configurations that are both symmetric and half-symmetric. But, being symmetric and half-symmetric uniquely determines a rectangle if we know a single vertex, so there is only one such configuration.

 


	\item (\textit{USAJMO 2024}) Let $a(n)$ be the sequence defined by $a(1)=2$ and $a(n+1)=(a(n))^{n+1}-1$ for each integer $n\geq1$. Suppose that $p>2$ is prime and $k$ is a positive integer. Prove that some term of the sequence $a(n)$ is divisible by $p^k$.


 



\textbf{Solution 1}

Lemma $1$:

 Given a prime $p$, a positive integer $k$, and an even $m$ such that $p^k|a(m)$, we must have that $p^{k+1}|a(m+2)$.

 
Proof of Lemma $1$:

 $a(m+1)\equiv (a(m))^{m+1}-1\equiv -1 \mod p^{k+1}$

 Then, $a(m+2)\equiv (a(m+1))^{m+2}-1\equiv (-1)^{m+2}-1\equiv 0 \mod p^{k+1}$

 
Therefore, by induction, if there exists an even integer $m$ such that $p|a(m)$, then for all integers $k$, $p^k|a(m+2k-2)$, so we are done if there exists an even $m$ such that $p|a(m)$.

 
Now, consider the case where there is some prime $p>2$ such that there are no even integers $m$ such that $p|a(m)$.

 Lemma $2$:

 In this case, we must have that $p|a(m)$ if $m\equiv -1 \mod p-1$ for all integers $m$.

 Proof of Lemma $2$:

 Suppose for the sake of contradiction that there exists some $m$ such that $m\equiv -1\mod p-1$ and $p$ does not divide $a(m)$. Then, we have $a(m+1)\equiv (a(m))^{m+1}-1\equiv 1-1\equiv 0\mod p$, by Fermat's Little Theorem. Since for all $p>2$, $p-1$ is even, then $m+1$ would be even. However this results in a contradiction.

 Then, we get that if $m\equiv -1\mod p-1$, then $0\equiv a(m)\equiv (a(m-1))^m-1\mod p\implies (a(m-1))^{m+1}\equiv 1\equiv a(m-1)\mod p$.

 Then, by LTE, $v_p(a(m))=v_p((a(m-1))^m-1)=v_p(a(m-1)-1)+v_p(m)>v_p(m)$. Since $\gcd(p-1,p)=1$, then $\gcd(p-1,p^k)=1$ for all positive integers $k$, so then by \href{https://artofproblemsolving.com/wiki/index.php/Chinese\_Remainder\_Theorem}{Chinese Remainder Theorem} there exists integers $m$ such that $m\equiv -1\mod p-1$ and $m\equiv 0\mod p^k$, so we are done $\square$

 
Remark: I think this is a very cool NT problem.

 
-bronzetruck2016

 


	\item (\textit{USAJMO 2024}) Let $n \ge 3$ be an integer. Rowan and Colin play a game on an $n \times n$ grid of squares, where each square is colored either red or blue. Rowan is allowed to permute the rows of the grid, and Colin is allowed to permute the columns of the grid. A grid coloring is $orderly$ if:


 
                \begin{itemize}
                    \item no matter how Rowan permutes the rows of the coloring, Colin can then permute the columns to restore the original grid coloring; and
                \end{itemize}


                
                \begin{itemize}
                    \item no matter how Colin permutes the columns of the coloring, Rowan can then permute the rows to restore the original grid coloring;
                \end{itemize}


                In terms of $n$, how many orderly colorings are there?


 



\textbf{Solution 1}

We focus on the leftmost column for simplicity. Let $m$ be the number of red squares in this column. We then have five cases:

 
1. $m=1$

 When Rowan permutes the rows of the coloring, we consider only the first column, which by the above contains $m=1$ red colors, so there are ${n \choose 1}=n$ ways to permute the first column’s rows. Thus every other column will have to contain one different permutation of the first column; otherwise, there will be at least one permutation of which there is no corresponding column.

 
Furthermore, each permutation will be different, so each row will contain one and only one red square, which also fulfills the case of if Colin permutes the coloring first. Thus there are $n\cdot (n-1)\cdot(n-2)\cdot\cdot\cdot2\cdot1=n!$ different colorings for this case (the same as choosing squares such as no square is in the same row or column as any other square).

 


 2. $m=n-1$

 This is essentially the same as case 1 except for the coloring; now there is one blue square and the rest are red squares. Thus there are also $n!$ different colorings for this case.

 


 3. $m=0$

 Since we have an entirely blue column, we are unable to have a column with $1$ red square only as doing so would leave one permutation that is not covered by at least one column (that space is being taken for the blank column). We are also unable to have a completely blue column as doing so would allow for Colin to shift the columns and in doing so fail for Rowan to shift back the columns. We also cannot have a column with any other number of red squares other than $0$ as will be shown below, so there is $1$ case here in which the entire coloring is red.

 


 4. $m=n$

 This is the same is an entire blue column, and, similar to above, we have $1$ coloring.

 


 5. $1<m<n-1$

 This is the final case and is equivalent to permuting for ${n \choose m}$ different ways. We must prove that this is greater than $n$ to show that the columns are not able to contain every possible permutation of this column for all values of $n$ such that $n>3$ (when $n=3$, there is no such positive integer $m$ that satisfies the conditions). Note that if we have any column with a different number of red squares, it is an unattainable column and is thus not optimal.

 
Lemma: Given that $m$ and $n$ are positive integers such that $1<m<n-1$ and $n>3$, it is true for all $m$ and $n$ that ${n \choose m}>n$.

 Proof: Assume that $m<\frac{n-1}{2}$.

 $\Leftrightarrow$ $m+1<n-m$

 $\Leftrightarrow$ $(m+1)!(n-m-1)!<m!(n-m)!$

 $\Leftrightarrow$ $\frac{n!}{m!(n-m)!}<\frac{n!}{(m+1)!(n-m-1)!}$

 $\Leftrightarrow$ ${n \choose m}<{n \choose m+1}$

 Similarly, we can prove that ${n \choose m}>{n \choose m+1}$ for $m>\frac{n-1}{2}$.

 Now we split our proof into two cases.

 Case 1: $n$ is even.

 The largest integer less than $\frac{n-1}{2}$ is $\frac{n}{2}-1$, so we know that:

 ${n \choose \frac{n}{2}}>{n \choose \frac{n}{2}-1}>\cdot\cdot\cdot>{n \choose 2}$

 by induction. On the other hand, the smallest integer greater than $\frac{n-1}{2}$ is $\frac{n}{2}$, so we know that:

 ${n \choose \frac{n}{2}}>{n \choose \frac{n}{2}+1}>\cdot\cdot\cdot>{n \choose n-2}$

 also by induction. Thus out of the given range for $m$ we know that ${n \choose 2}$ and ${n \choose n-2}$ are the minimum values, and all that is left is to prove that they are both greater than $n$. Furthermore, since ${n \choose 2}={n \choose n-2}$, we only have to prove that ${n \choose 2}>n$.

 We start with the given: $n>3$

 $\Leftrightarrow$ $\frac{n-1}{2}>1$

 $\Leftrightarrow$ $\frac{n(n-1)}{2}>n$

 $\Leftrightarrow$ $\frac{n!}{2!(n-2)!}>n$

 $\Leftrightarrow$ ${n \choose 2}>n$

 Thus we have proven the inequality for all even $n$.

 Case 2: $n$ is odd.

 The greatest integer less than $\frac{n-1}{2}$ is $\frac{n-3}{2}$, so we know that:

 ${n \choose \frac{n-1}{2}}>{n \choose \frac{n-3}{2}}>\cdot\cdot\cdot>{n \choose 2}$

 by induction. On the other hand, the smallest integer greater than $\frac{n-1}{2}$ is $\frac{n+1}{2}$, so we know that:

 ${n \choose \frac{n+1}{2}}>{n \choose \frac{n+3}{2}}>\cdot\cdot\cdot>{n \choose n-2}$

 also by induction. Since ${n \choose \frac{n+1}{2}}={n \choose \frac{n-1}{2}}$, we know that once again, ${n \choose n-2}={n \choose 2}$ is the minimum of the given range for $m$, and the same proof applies. Thus, the inequality holds true for odd and in turn all positive integers $n>3$.

 
As a result, due to our lemma, there are always more permutations of the columns than the number of columns itself, so there will always exist a permutation of the column such that there are no corresponding original columns of which to match with. Thus there are no solutions for this case.

 
In conclusion, there are a total of $2\cdot n!+2$ different colorings for which the above apply.

 



\textbf{Solution 2}

Lemma 1: Each row and column must have the same number of red squares.

 Proof: Suppose two rows do not have the same number of red squares. Suppose Rowan permutes the two rows onto each other. Then, because the two rows have a different number of red squares, there is no way for Colin to permute the columns, which permutes the squares in the row, to match the other row because the two rows have a different number of red squares. By a similar proof, the columns must have the same number of red squares as each other. Assume that each row has $m$ red squares. Then, there are $nm$ red squares in total. Because each column must have the same number of red squares as the other columns, the columns also must have $m$ red squares each.

 
Now, we want to find which values of $m$ can produce a valid orderly coloring.

 
Lemma: $C(n,k)>n$ for $k\neq 0,1,n-1,n.$ This is needed to prove the next lemma and is the most mathematical part of the solution.

 Proof: Start with $C(n,1)=n.$ Each time we increase $k$ by $1$ until $k=\lfloor0.5n\rfloor,$ we are multiplying the previous value of $C(n,k)$ by a fraction more than $1.\text{*}$ Thus, $C(n,k)>n$ for $2\le n \le \lfloor0.5n\rfloor.$ Because $C(n,k)=C(n,n-k),$ the lemma is proven for the rest of the conjectured values. 

 
Lemma 2: If $m\neq 0,1,n-1,n,$ the coloring is not orderly.

 Proof: Let's take a look at a particular column, where assume there are $2\le m\le n-2$ red squares. When Rowan permutes the rows, the squares in this column have $C(n,m)$ possible permutations. In order for Colin to be able permute the columns to return the grid to its original state, all of these $C(n,m)$ permutations must be columns of the original grid. However, because there are only $n$ columns. By the previous lemma, there are more permutations than columns, so the we are done.

 
Lemma 3: If a permutation satisfies $m=0,1,n-1,n$ and satisfies lemma $1,$ it is orderly.

 Proof: The $m=0,n$ case is trivial: the grid never changes! For $m=1,n-1,$ we only have to prove the $n=1$ case because of symmetry: we can change red to blue and blue to red and it won't change the orderliness of the coloring as the overall arrangement of the different colors won't be changed. For the $m=1 case,$ when Rowan permutes the rows, all of the columns will still be different columns of the original grid (as each column must have one red square, the rows of the red squares will be different) and can be permuted back to the original arrangement by Colin. Thus, this lemma is true.

 
Finally, there is $1$ coloring for the $m=0$ case, and $n!$ ways to permute the red squares to see which column values are assigned to the rows in the $m=1$ case. There are similarly $n!$ colorings for the $m=n-1$ case, and $1$ coloring for the $m=0$ case. Thus, the answer is $\boxed{2n!+2}.$

 


 $\text{*}$ Specifically, $\frac{C(n,k)}{C(n,k-1)}=\frac{n-k+1}{k-1}$ which is greater than $1$ for $2\le n \le \lfloor0.5n\rfloor.$

 \textbf{Another way to prove that $\binom{n}{m}>n$} Suppose we have $2\le{m}\le\lfloor{\frac{n}{2}}\rfloor$. Our goal is to show that

 
 \[\binom{n}{m}>n\]

 
 \[\frac{n}{m}\cdot\frac{n-1}{m-1}\cdot\frac{n-2}{m-2}\cdots\frac{n-m+1}{1}>n\]

 
 \[\frac{1}{m}\cdot\frac{n-1}{m-1}\cdot\frac{n-2}{m-2}\cdots\frac{n-m+1}{1}>1\]

 
 \[\frac{n-1}{m-1}\cdot\frac{n-2}{m-2}\cdots\frac{n-m+1}{m}>1\].

 Each of the above fractions is greater than one because $2\le{m}\le\lfloor{\frac{n}{2}}\rfloor$, so we are done.

 


	\item (\textit{USAJMO 2024}) Find all functions $f:\mathbb{R}\rightarrow\mathbb{R}$ that satisfy

 \[f(x^2-y)+2yf(x)=f(f(x))+f(y)\]
for all $x,y\in\mathbb{R}$.


 



\textbf{Solution 1}

We claim the solutions are $\boxed{f\equiv 0, x^2, -x^2}$. These solutions clearly work. Let the given assertion be $P(x, y)$. $P(0, 0)$ yields $f(f(0)) = 0$. Then, $P(f(0), y)$ produces 
 \[f(0)+f(y)=f(f(0)^2-y)\] Letting $y = 0, f(0)^2$ gives $2f(0)=f(f(0)^2)$ and $f(f(0)^2) = 0$, so $f(0)= 0$. Now, $P(x, x^2)$ and $f(x, 0)$ implies 
 \[x^2f(x)=f(f(x))=f(x^2) \tag{1}\]Notice that this implies $f$ is even, i.e. $f(x)=f(-x)$. Furthermore, 
 \[f(x_1)=f(x_2)\implies f(f(x_1))=f(f(x_2)) \implies x_1^2 f(x_1) = x_2^2 f(x_2)\]
 \[\implies x_1^2 = x_2^2 \text{ or } f(x_1)=f(x_2) = 0 \tag{2}\] To rule out the second case, suppose $f(c)=0, c\neq 0$. Assuming $f$ is not identically 0, let $f(y) \neq 0$, so $P(c, y)$ and $P(c, -y)$ gives
 \[f(c^2-y)=f(y)=f(-y) = f(c^2+y) \implies (c^2-y)^2 = (c^2+y)^2 \implies y = 0\] a contradiction. Thus, by $(1)$ and $(2)$, we conclude $|f(x)| = x^2$. 

 Lastly, to eliminate pointwise trap, assume $f(a)=a^2, f(b)=-b^2, a, b \neq  0$. Since $f$ is even, assume $b< 0$. Then, $P(a, b)$ yields 
 \[f(a^2-b)+2a^2b=f(a^2)-b^2\] Casework on the signs of $f(a^2-b)$ and $f(a^2)$ yields that either $a$ or $b$ is 0, a contradiction. Thus, the only solutions are $\boxed{f\equiv 0, x^2, -x^2}$.

 -sillybone

 



\textbf{Solution 2}

Plugging in $y$ as $0:$\begin{equation}f(x^2)=f(f(x))+f(0) \text{ } (1)\end{equation}Plugging in $x, y$ as $0:$
 \[f(0)=f(f(0))+f(0)\]or
 \[f(f(0))=0\]Plugging in $x$ as $0:$
 \[f(-y)+2yf(0)=f(f(0))+f(y),\]but since $f(f(0))=0,$\begin{equation}f(-y)+2yf(0)=f(y) \text{ } (2)\end{equation}Plugging in $y^2$ instead of $y$ in the given equation:
 \[f(x^2-y^2)+2y^2f(x)=f(f(x))+f(y^2)\]Replacing $y$ and $x$:
 \[f(y^2-x^2)+2x^2f(y)=f(f(y))+f(x^2)\]The difference would be:\begin{equation}f(x^2-y^2)-f(y^2-x^2)+2y^2f(x)-2x^2f(y)=f(f(x))-f(x^2)-(f(f(y))-f(y^2)) \text{ } (3)\end{equation}The right-hand side would be $f(0)-f(0)=0$ by $(1).$ Also,
 \[f(x^2-y^2)-f(y^2-x^2)=2(x^2-y^2)f(0)\] by $(2)$So, $(3)$ is reduced to:
 \[2(x^2-y^2)f(0)+2y^2f(x)-2x^2f(y)=0\]Regrouping and dividing by 2:
 \[y^2(f(x)-f(0))=x^2(f(y)-f(0))\]
 \[\frac{f(x)-f(0)}{x^2}=\frac{f(y)-f(0)}{y^2}\]Because this holds for all x and y, $\frac{f(x)-f(0)}{x^2}$ is a constant. So, $f(x)=cx^2+f(0)$.This function must be even, so $f(y)-f(-y)=0$.So, along with $(2)$, $2yf(0)=0$ for all $y$, so $f(0)=0$, and $f(x)=cx^2$.Plugging in $cx^2$ for $f(x)$ in the original equation, we get:
 \[c(x^4-2x^2y+y^2)+2cx^2y=c^3x^4+cy^2\]
 \[c(x^4+y^2)=c(c^2x^4+y^2)\]So, $c=0$ or $c^2=1.$All of these solutions work, so the solutions are $f(x)=-x^2, 0, x^2$.

 -codemaster11

 



\textbf{Solution 3}

Let our equation be $P(x,y)$. We start by plugging in some initial values:

 $y=x^2:\; f(0)+2x^2f(x) = f(f(x))+f(x^2) \;\;\;\; (1)$

 $y=0:\; f(x^2) = f(f(x))+f(0) \;\;\;\; (2)$

 $x=0:\; f(-y) + 2yf(0) = f(f(0)) + f(y) \;\;\;\; (3)$

 Plugging in $x=1$ into $(3)$ gives
 \[f(1) = f(f(1)) + f(0) \;\;\;\; (4).\]From $(1)$, we get
 \[f(0) + 2x^2f(x) = 2f(x^2)-f(0) \implies x^2f(x)+f(0) = f(x^2)\]Substituting in what we have in $(3)$ gives
 \[x^2f(x)+f(0) = f(0) = f(f(x)) \implies x^2f(x) = f(f(x)).\]Plugging in $x=1$ gives
 \[f(1)=f(f(1)) \;\;\;\; (5).\]As a result, $(4)$ becomes $f(0)=0$.

 Now, $(3)$ becomes
 \[f(x^2) = f(f(x)) \;\;\;\; (6)\]and $(2)$ becomes
 \[f(y)=f(-y) \;\;\;\; (7).\]Note that $f\equiv 0$ is a solution. Now, assume $f(x) \neq 0$.

 \textbf{Claim:} $f$ is injective over $\mathbb{R}^{+}$.

 Let $f(a) = f(b) \neq 0$ with $a,b>0$. Plugging in $x=a, y=b^2$ and $x=b, y=a^2$ into $P$ gives us
 \[f(a^2-b^2)+2b^2f(a) = f(a^2)+f(b^2)\]
 \[f(b^2-a^2)+2a^2f(b) = f(b^2)+f(a^2)\]Subtracting, and using $(7)$ gives us $2(a^2-b^2)f(a) = 0$, which implies that either $f(a)=0$ or $a=\pm b$. Either way leads to contradiction. Thus, $f$ is injective. $\square$

 As a result, $(6)$ becomes $f(x)=\pm x^2$. Piecing everything yields $f(x) = 0, \pm x^2$.

 It just remains to verify these solutions work, and doing so is quite trivial;
 \[f(x)=0:\; 0+0 = 0+0,\]
 \[f(x)=x^2:\; (x^2-y)^2+2yx^2 = x^4+y^2,\]
 \[f(x)=-x^2:\; -(x^2-y)^2-2yx^2 = -x^4-y^2,\]all of which are obviously true.

 



\textbf{Solution 4}

Start with $y = 0$ to get: $f(x^{2}) = f(f(x)) + f(0)$. Then substitute $x = 0$ to get: $f(-y) + 2yf(0) = f(f(0)) + f(y)$. We will show that $f(f(0)) = 0$. Substitute $x = y = 0$ to get $f(0) = f(f(0)) + f(0)$ which yields the desired $f(f(0)) = 0$. Hence, $f(-x) + 2xf(0) = f(x)$ since $y$ is just a dummy variable. Now we solve for $f(0)$ in our two equations we found. We have that $f(0) = f(x^{2}) - f(f(x))$ from the first equation and that $f(0) = \frac{f(x) - f(-x)}{2x}$ from the second equation.

 Hence, we set these equal to get:

 $2xf(x^{2}) - 2xf(f(x)) = f(x) - f(-x)$

 Replace $x$ with $-x$ to get:

 $-2xf(x^{2}) + 2xf(f(-x)) = f(-x) + f(x)$

 Now add to get:

 $2x(f(f(-x)) - f(f(x))) = 0$

 Now since $x$ can't obviously be equal to $0$, we need $f(f(-x)) = f(f(x))$.

 Case 1: $f$ is an injective function

 Hence, $f(-x) = f(x)$ and thus $f(x)$ is an even function.

 \textbf{Claim:} $f$ cannot have degree at least $4$

 \textbf{Proof:}

 Assume $f$ is a polynomial. Assume there exists a solution of which $f$ has degree at least $4$. Hence, $f(x) = a_n x^{2n} + a_{n - 1} x^{2n - 2} + ... + a_1 x^{2} + a_0$ where $n \geq 2$. If we go back to the original equation in the problem statement and plug in $y = x^{2}$, then we have: 

 $f(0) + 2x^{2} f(x) = f(f(x)) + f(x^{2})$

 So we have:

 $a_0 + 2x^{2}f(x) = a_n f(x)^{2n} + ... + a_1 f(x)^{2} + a_0$

 Hence, $2x^{2} f(x) = a_n f(x)^{2n} + ... + a_1 f(x)^{2}$

 But we know $f$ has degree $2n$ so the $LHS$ will have degree $2 + 2n$ but the $RHS$ will have degree $(2n)(2n) = 4n^{2}$. Solving $4n^{2} = 2n + 2$ will have no integer solutions. Recall that $n \geq 2$. Clearly, no functions will satisfy this equation and the claim is proven. $\square$

 We have shown that $f(x)$ can have degree at most $2$ if it is an injective function. Hence, let $f(x) = ax^{2} + bx + c$. Recall that we had shown earlier that $f(f(0)) = 0$. Hence, $a f(0)^{2} + b f(0) + c = 0$ but we know $f(0) = c$. Hence, $ac^{2} + bc + c = 0$ which yields $c(ac + b + 1) = 0$. Clearly, $c = 0$ has to occur. Hence, $f(x) = ax^{2} + bx$. Now, recall we showed that $f(x^{2}) = f(f(x)) + f(0)$. Hence, $ax^{4} + bx^{2} = af(x)^{2} + bf(x) + f(0)$. Plugging and simplifying, we eventually get:

 $x^{4} (a - a^{3}) - 2a^{2} b x^{3} - abx^{2} (b + 1) = f(0) = constant$

 Hence, all the $x$ terms must cancel out for that to be a constant. This can be achieved by either $a = 0$ and $b$ can be anything OR $b = 0$ and $a - a^{3} = 0$. Now if $a = 0$, then $f(x) = bx$ but recall that $f(x)$ is an even function and hence $b = 0$ is the only suitable option for a linear function to turn into an even function. Hence, $f(x) = 0$ is a solution. If $b = 0$ and $a - a^{3} = 0$, then $a = 1, -1, 0$ and hence $f(x) = x^{2}, -x^{2}$ are two more new solutions.

 Now we have to consider if $f(x)$ is not a polynomial. 

 \textbf{Claim:} $f$ has to be a polynomial to satisfy the given problem statement

 \textbf{Proof:}

 We will first show that $f(f(x)) = x^{2} f(x)$. Substitute $y = \frac{x^{2}}{2}$:

 $f(\frac{x^{2}}{2}) + x^{2} f(x) = f(f(x)) + f(\frac{x^{2}}{2})$. Hence, the lemma. $\square$

 We have already shown that if $f$ is injective, then $f$ is even and this proof lies in the case of when $f$ is injective. Hence, $f(x)$ is an even function. Recall that $f(x^{2}) = f(f(x)) + f(0)$. Hence, $f(x^{2}) = x^{2} f(x) + f(0)$. Now we plug in $x = 1$ to get $f(1) = f(1) + f(0)$ and hence $f(0) = 0$. In other words, $f(x^{2}) = x^{2} f(x)$. (Haven't finished proof yet)

 Case 2: $f$ is not an injective function

 We can prove that $f$ has to be injective. This is shown in sml1809's solution above in which we claim $f(a) = f(b) \neq 0$ and substitute $x = a, y = b^{2}$ and $x = b, y = a^{2}$ to prove the claim that $f$ has to be injective.

 



\textbf{Solution 5(Another solution)}

We will show $x^2 f(x) = f(f(x))$. First, the original equation is

 $f(x^2 - y) + 2yf(x) = f(f(x)) + f(y)$

 We will make the substitution $y = x^2 - 1$

 $f(1) + 2(x^2 - 1)f(x) = f(f(x)) + f(x^2 - 1)$

 We will now take the original equation again and substitute $y = 1$

 $f(x^2 - 1) + 2f(x) = f(f(x)) + f(1) \implies f(x^2 - 1) = f(f(x)) + f(1) - 2f(x)$

 Now we will substitute this new representation of $f(x^2 - 1)$ into our equation to get

 $f(1) + (2x^2 - 2)f(x) = f(f(x) + f(f(x)) + f(1) - 2f(x) \implies (2x^2)f(x) = 2f(f(x)) \implies x^2 f(x) = f(f(x)) \square.$

 We will now plug this in to our given equation

 $f(x^2 - y) + 2yf(x) = f(f(x)) + f(y) \implies f(x^2 - y) + 2yf(x) = x^2 f(x) + f(y) \implies f(x^2 - y) - f(y) = f(x)(x^2 - 2y)$. 

 Now we will make the substitution $x = y$

 $f(x^2 - x) - f(x) = f(x)(x^2 - 2x) \implies f(x^2 - x) = f(x)(x^2 - 2x + 1) \implies f(x(x - 1)) = f(x)(x - 1)^2$. This immediately raises a conjecture. If we have two real numbers $a, b$ where $f$ is defined at these two real numbers, then

 $f(ab) = f(a) b^2$ 

 Let's prove this. From our equation $f(x^2 - y) + 2yf(x) = x^2 f(x) + f(y)$, we can substitute $y = -x$ to get

 $f(x^2 + x) - 2xf(x) = x^2 f(x) + f(-x) \implies f(x^2 + x) = f(x)(x^2 + 2x) + f(-x)$. If we can show $f(-x) = f(x)$, then $f(x^2 + x) = f(x(x + 1)) = f(x)(x^2 + 2x + 1) = f(x) (x + 1)^2 \implies f(x(x + 1) = f(x) (x + 1)^2$ and we should be done as we can now prove $f(ab) = f(a) b^2$ by induction. So let's prove $f(-x) = f(x)$. 

 The proof is simple. From our given equation $f(x^2 - y) + 2yf(x) = f(f(x)) + f(y)$, let's substitute $x = 0$ to get

 $f(-y) = f(f(0)) + f(y)$. 

 We will prove $f(0) = 0$. From our given equation $f(x^2 - y) + 2yf(x) = f(f(x)) + f(y)$, we will plug in $x = y = 1$ to get

 $f(0) + 2f(1) = f(f(1)) + f(1) \implies f(0) + f(1) = f(f(1)) \implies f(0) + f(1) = 1^2 f(1) \implies f(0) = 0$ following $f(f(x)) = x^2 f(x)$. Henceforth, we go back to $f(-y) = f(f(0)) + f(y) \implies f(-y) = f(0) + f(y) \implies f(-y) = f(y) \implies f(-x) = f(x) \square.$

 Hence, as $f(-x) = f(x)$, we had concluded before that if this was the case, 

 $f(x(x + 1)) = f(x) (x + 1)^2$ and $f(x(x - 1)) = f(x) (x - 1)^2$ where the second equation we had derived before proving $f(-x) = f(x)$. From these two equations, we can generate any two real numbers that differ by $1$ where

 $f(ab) = f(a) b^2$ where $a - b = 1$. 

 We will show this happens for all real $a, b$ as we have proved only part of our conjecture. Note that $f(a) = f(-a)$ has been proved. 

 We will assume $f(ab) = f(a) b^2$ only holds true if and only if $a - b = 1$. Now let $b = -1$. We have

 $f(-a) = f(a)$ which holds true for all real numbers. However, we assumed the above was true if and only if $a - b = 1$. This would imply the above was true if and only if $a = 0$ which is not the case. Hence, this is a contradiction and we have proved our conjecture. $\square.$

 Now, we have shown $f(ab) = f(a) b^2$ is true for all real $a, b$. Let $a = 1$ to get

 $f(b) = b^2$ which means $f(x) = x^2$ is a solution. 

 Now consider $ab = x \implies b = \frac{x}{a}$. Now $f(x) = f(a) \frac{x^2}{a^2} \implies \frac{f(x)}{x^2} = \frac{f(a)}{a^2}$. From here we can generate many quadratic functions. Knowing $f(x) = x^2$ is a solution gives us motivation for $f(x) = -x^2$ which we can confirm will satisfy the given equation. We can prove this also

 Given the original equation $f(x^2 - y) + 2yf(x) = f(f(x)) + f(y)$, we have found $f(x) = x^2$ is a solution. We will show that if $f(x)$ is a solution, so is $-f(x)$. First, we'll assume something different, then prove by contradiction. We claim that if $f(x)$ is a solution, $-f(x)$ is not a solution. Let's plug in $f(x) = x^2$ since we already know that's a solution. We have

 $(x^2 - y)^2 + 2x^2 y = x^4 + y^2 \implies x^4 - 2x^2 y + y^2 + 2x^2 y = x^4 + y^2$ which holds. 

 So now since we know $f(x) = x^2$ is a solution, we had claimed that $-f(x) = -x^2$ was not a solution. Hence we can plug this in to get

 $-(x^2 - y)^2 - 2x^2 y = f(-x^2) - y^2 \implies -(x^2 - y)^2 - 2x^2 y = -x^4 - y^2$ of which this is exactly the negation of the equation we got for $f(x) = x^2$. Therefore $-f(x)$ is indeed a solution and we have a contradiction with $-f(x)$ not being a solution. Thus $-f(x)$ has to be a solution. $\square$.

 Now assuming $f$ is constant, we have

 $f(x) = c \implies c + 2yc = x^2 c + c \implies 2yc = x^2 c \implies c = 0$ is a solution. Hence $f(x) = 0$ is also a solution and so $f(x) = x^2, -x^2, 0$

 



\textbf{Solution 6}

Denote $P(x,y)$ as the problem statement's assertion. Denoting $a=f(0)$, We have the following:\begin{align}P(x,0) \implies& f(x^2)=f(f(x))+a & (1) \\P(0,x) \implies & f(-x)+2ax=f(a)+f(x)&(2)\\P\left(x,\frac{1}{2}x^2\right) \implies & x^2f(x)=f(f(x))&(3)\\P(x,y^2) \implies& f(x^2-y^2)+2y^2f(x)=f(f(x))+f(y^2)&(4)\\P(y,x^2) \implies & f(y^2-x^2)+2x^2f(y)=f(f(y))+f(x^2)&(5)\end{align}Combining $(1)$ and $(3)$, we have $f(x^2)=x^2f(x)+a$. Plugging in $x=1$, we have $f(1)=f(1)+a \implies f(0)=0$. Rewriting $(1)$ and combining with $(3)$ using this fact, we have 
 \[f(x^2)=f(f(x))=x^2f(x).\]Rewriting $(2)$ using this fact, we have $f(-x)=f(x)$ (i.e., $f$ is even). We may henceforth limit $x,y$ to the positive reals without loss of generality. Rewriting $(4)$ and $(5)$ and employing the fact that $f$ is even, we have \begin{align*}f(x^2-y^2)+2y^2f(x)&=x^2f(x)+y^2f(y)\\f(y^2-x^2)+2x^2f(y)&=y^2f(y)+x^2f(x)\\\implies f(x^2-y^2)+2y^2f(x)&=f(y^2-x^2)+2x^2f(y)\\\implies 2y^2f(x)&=2x^2f(y)\\\implies \frac{f(x)}{x^2}&=\frac{f(y)}{y^2}\end{align*}
Since the choice of $x,y$ was arbitrary, we find that $f(x)/x^2$ is constant for all $x$, so $f(x)=cx^2$ for some constant $c$. Plugging into the original functional equation, we have \begin{align*}c(x^2-y)^2+2y(cx^2)&=c(cx^2)^2+cy^2\\\iff cx^4-2cx^2y+cy^2+2cx^2y&=c^3x^4+cy^2\\\iff (c-c^3)x^4&=0\end{align*}
for all $x,y$. Thus, it must be that $c\in \{-1,0,1\}$, so the only valid solutions are $\boxed{f(x)=0, \pm x^2}$. $\blacksquare$

 


	\item (\textit{USAJMO 2024}) Point $D$ is selected inside acute triangle $ABC$ so that $\angle DAC=\angle ACB$ and $\angle BDC=90^\circ+\angle BAC$. Point $E$ is chosen on ray $BD$ so that $AE=EC$. Let $M$ be the midpoint of $BC$. Show that line $AB$ is tangent to the circumcircle of triangle $BEM$.


 



\textbf{Solution 1}

Let $\angle DBT = \alpha$ and $\angle BEM = \beta$.  Extend AD intersects BC at point T, then TC = TA, TE is perpendicular to AC

 Thus, AB is the tangent of the circle BEM

 Then the question is equivalent as  the $\angle ABT$ is the auxillary angle of $\angle BEM$.

 \textit{ontinued}

 



\textbf{Solution 2}

We notice that $\angle DAC = \angle ACB$, which implies symmetry about the perpendicular bisector of $\overline{AC}$. Therefore, we can extend $\overline{AD}$ to intersect the circumcircle at $Q$, and by symmetry about the perpendicular bisector, $ABQC$ is an isoceles trapezoid. Now, we realize that the midpoint condition is annoying, which really suggests that to simplify things, we take a homothety at $B$ sending $M$ to $C$, and we just need to prove that the circumcircle of $\triangle BE'C$ is tangent to line $\overline{AB}$, where $E'$ is the resultant point from a homothety at $E$. However, funnily enough, $E$ lies on the perpendicular bisector of $\overline{AC}$ because $\overline{AE}=\overline{EC}$, and the perpendicular bisector gets sent to an altitude from $Q$ under homothety, which means $E'$ is simply the intersection of $\overline{BD}$ and the altitude from $Q$ onto $\overline{BC}$! Call that point $P$. We can now ignore points $E$ and $M$. Note that we have not used the $\angle BDC$ condition yet, which means we should probably try angle chasing now. Call $\angle BAC=a, \angle ABC = b$. By angle condition, $\angle BDC = 90+a, \angle CDP = 90-a$. Now, $\angle QCA = \angle BAC = a$, so $\angle PQC = 90 - \angle ACQ = 90-a$. A miracle occurs! $\angle CDP = \angle ACQ$, so quadrilateral $QDPC$ is cyclic! $\angle ABC = \angle AQC = b$ from isoceles trapezoid, and from cyclicity $\angle DPC = 180-b = \angle BPC$. Now we have $\angle ABC = b$ and $\angle BPC = 180-b$, which looks a lot like tangency condition. We realize that it simply is the tangency condition if we extend line $\overline{AB}$ to get the supplementary angle, so we are done.

 


\end{enumerate}